\documentclass[10pt,a4paper]{article}

\usepackage[utf8]{inputenc}
\usepackage[T1]{fontenc}
\usepackage[french]{babel}
\usepackage{amsmath,amssymb,amsthm}
\usepackage{geometry}
\usepackage{booktabs}
\usepackage{enumitem}
\usepackage{dsfont}
\usepackage{tikz}
\usepackage{hyperref}

\geometry{hmargin=1.5cm, top=1cm, bottom=1.2cm, footskip=0.5cm}
\setlength{\parskip}{3pt}
\setlength{\parindent}{0pt}
\setlist{itemsep=1pt,topsep=2pt}

\newcommand{\A}{\mathcal{A}}
\newcommand{\R}{\mathcal{R}}
\newcommand{\N}{\mathcal{N}}

\title{Projet SDP -- Apprentissage de modèles MR-Sort}
\author{Côme-Alexis Puech \and Hugo Vigna \and Salomé Fonvielle}
\date{Octobre 2026}

\begin{document}

\maketitle

Ce rapport présente notre programme pour le Jalon 1, qui apprend les paramètres d'un modèle MR-Sort à partir
d'un learning set en résolvant un MILP avec Gurobi. La première partie construit ce MILP, la seconde décrit son
implémentation et sa validation. Les réponses détaillées aux questions préliminaires i à iv, sur l'exemple de
l'admission d'étudiants, sont en annexe.

\section*{Formulation du MILP}

Cette partie reprend le raisonnement des questions iv à viii bis du sujet.

\subsection*{Maximiser une marge}

Avec deux catégories, on note $A^*$ et $R^*$ les objets acceptés et rejetés du learning set, et $s(x)=\sum_{i:\,x_i\geq b_i} w_i$
le score de $x$, c'est-à-dire le poids des critères sur lesquels il atteint la frontière $b$. Le modèle classe bien $x$ si
$s(x)\geq\lambda$ pour $x\in A^*$ et $s(x)<\lambda$ pour $x\in R^*$. On cherche donc $w$ et $\lambda$ qui maximisent une marge
$\alpha$, le plus petit écart entre un score et $\lambda$ :
\[
\max\ \alpha \quad\text{s.c.}\quad s(x)\geq\lambda+\alpha\ \ (x\in A^*),\qquad s(x)\leq\lambda-\alpha\ \ (x\in R^*),\qquad \textstyle\sum_{i\in\N} w_i=1 .
\]
Si $\alpha^*>0$, tous les objets sont bien classés, sinon aucun modèle ne colle parfaitement aux données.
Quand $b$ est connu, $s(x)$ est linéaire en $w$ et c'est un simple PL. Pour apprendre aussi $b$, il faut exprimer
$s(x)$ de façon linéaire en fonction de $b$.

\subsection*{Linéarisation pour apprendre les frontières}

\textbf{Variables $\delta_i(x)$.} Pour savoir si $x$ valide le critère $i$, on introduit une variable binaire $\delta_i(x)=\mathds{1}_{\{x_i\geq b_i\}}$. Comme $b_i$ est
une variable, cette indicatrice n'est pas linéaire. On l'impose avec deux contraintes, une par sens de l'équivalence
($\forall i\in\N$, $\forall x$) :
\begin{align*}
x_i-b_i &\geq M\,(\delta_i(x)-1) &\qquad (1)\quad \delta_i(x)=1 \Rightarrow x_i\geq b_i \\
x_i-b_i &\leq M\,\delta_i(x)-\varepsilon &\qquad (2)\quad \delta_i(x)=0 \Rightarrow x_i< b_i
\end{align*}
On normalise les données dans $[0,1]$ sur chaque critère, ce qui ne change pas les affectations. On a alors
$-1\leq x_i-b_i\leq 1$, donc avec $M=2$ la contrainte qui ne correspond pas à la valeur de $\delta_i(x)$ est toujours
vérifiée. Comme le solveur ne gère pas les inégalités strictes, on écrit $x_i<b_i$ sous la forme $x_i\leq b_i-\varepsilon$,
avec $\varepsilon$ plus petit que l'écart entre deux valeurs voisines des données.

\begin{center}
\begin{tabular}{cccc}
\toprule
 & Contrainte (1) & Contrainte (2) & Conclusion \\
\midrule
$\delta_i(x)=1$ & $x_i-b_i\geq 0$ & $x_i-b_i\leq M-\varepsilon$ \ (toujours vraie) & $x_i\geq b_i$ \\
$\delta_i(x)=0$ & $x_i-b_i\geq -M$ \ (toujours vraie) & $x_i-b_i\leq-\varepsilon$ & $x_i<b_i$ \\
\bottomrule
\end{tabular}
\end{center}

\textbf{Variables $w_i(x)$.} Le score s'écrit alors $s(x)=\sum_{i\in\N} w_i\,\delta_i(x)$, mais $w_i\,\delta_i(x)$ est un produit
de deux variables. On le remplace par une variable $w_i(x)\geq 0$ et les contraintes
\begin{align*}
w_i \geq w_i(x) &\geq 0 && (3)\\
\delta_i(x) \geq w_i(x) &\geq \delta_i(x)+w_i-1 && (4)
\end{align*}
qui forcent $w_i(x)=w_i\,\delta_i(x)$ (les autres bornes sont toujours vraies car $0\leq w_i\leq 1$) :

\begin{center}
\begin{tabular}{cccc}
\toprule
 & Borne inférieure & Borne supérieure & Conclusion \\
\midrule
$\delta_i(x)=1$ & $w_i(x)\geq w_i$ \ (4) & $w_i(x)\leq w_i$ \ (3) & $w_i(x)=w_i$ \\
$\delta_i(x)=0$ & $w_i(x)\geq 0$ \ (3) & $w_i(x)\leq 0$ \ (4) & $w_i(x)=0$ \\
\bottomrule
\end{tabular}
\end{center}

Le score $s(x)=\sum_{i\in\N} w_i(x)$ est donc linéaire, et on n'a plus besoin de connaître $b$ à l'avance.

\subsection*{Le MILP à deux catégories}

On reprend le programme de la marge avec ce score et on ajoute les contraintes (1) à (4) ($M=2$, $\varepsilon$ petit).

\medskip
\textbf{Variables.}
\begin{itemize}
  \item du modèle : $w_i\geq 0$, $b_i\in[0,1+\varepsilon]$ pour $i\in\N$ (au-delà de 1, aucun objet n'atteint la frontière), et $\lambda\in[0,1]$ ;
  \item de la marge : $\alpha$ ;
  \item pour chaque objet $x$ et chaque critère $i$ : $\delta_i(x)\in\{0,1\}$ et $w_i(x)\geq 0$.
\end{itemize}

\medskip
\textbf{Programme.}
\begin{align*}
\max\quad & \alpha \\[4pt]
\text{s.c.}\quad & \textstyle\sum_{i\in\N} w_i(x) \geq \lambda + \alpha && \forall x\in A^* && \text{(acceptés)}\\
& \textstyle\sum_{i\in\N} w_i(x) \leq \lambda - \alpha && \forall x\in R^* && \text{(rejetés)}\\[4pt]
& x_i-b_i \geq M\,(\delta_i(x)-1) && \forall i,\ \forall x && (1)\\
& x_i-b_i \leq M\,\delta_i(x)-\varepsilon && \forall i,\ \forall x && (2)\\[4pt]
& w_i \geq w_i(x) \geq 0 && \forall i,\ \forall x && (3)\\
& \delta_i(x) \geq w_i(x) \geq \delta_i(x)+w_i-1 && \forall i,\ \forall x && (4)\\[4pt]
& \textstyle\sum_{i\in\N} w_i = 1
\end{align*}

\medskip
Les $\delta_i(x)$ sont binaires, c'est donc un programme linéaire en nombres entiers (MILP) avec $n\,|A^*\cup R^*|$
variables binaires. Comme plus haut, si $\alpha^*>0$ le modèle appris restitue tout le learning set.
Si les données ne sont pas séparables ($\alpha^*<0$), maximiser $\alpha$ réduit la pire erreur, ce qui n'est pas la même chose
que maximiser le nombre d'objets bien classés. Ce cas ne se produit pas dans nos tests, puisque les learning sets sont
générés par un modèle MR-Sort.

\subsection*{Généralisation à $p$ catégories}

Les catégories vont de $\mathcal{C}^1$ (la pire) à $\mathcal{C}^p$ (la meilleure) et on cherche $p-1$ frontières
$b^1,\dots,b^{p-1}$, $b^h$ séparant $\mathcal{C}^h$ et $\mathcal{C}^{h+1}$. Les poids et $\lambda$ sont les mêmes pour toutes les
frontières, seules les valeurs $b^h_i$ changent. Un objet est dans $\mathcal{C}^h$ s'il franchit $b^{h-1}$ mais pas $b^h$.
On reprend donc les variables $\delta^h_i(x)$, $w^h_i(x)$ et les contraintes (1) à (4) pour chaque frontière $h$,
avec le score $s_h(x)=\sum_{i\in\N} w^h_i(x)$.

\medskip
\textbf{Programme.}
\begin{align*}
\max\quad & \alpha \\[4pt]
\text{s.c.}\quad & s_{h-1}(x)=\textstyle\sum_{i\in\N} w^{h-1}_i(x) \geq \lambda + \alpha && \forall x\in\mathcal{C}^h,\ h\geq 2 && \text{(franchit $b^{h-1}$)}\\
& s_h(x)=\textstyle\sum_{i\in\N} w^h_i(x) \leq \lambda - \alpha && \forall x\in\mathcal{C}^h,\ h\leq p-1 && \text{(ne franchit pas $b^h$)}\\[4pt]
& \text{(1) à (4) avec } b^h,\ \delta^h_i(x),\ w^h_i(x) && \forall i,\ \forall x,\ \forall h \\[4pt]
& b^{h-1}_i \leq b^h_i && \forall i,\ 2\leq h\leq p-1 && \text{(frontières ordonnées)}\\
& \textstyle\sum_{i\in\N} w_i = 1
\end{align*}

\medskip
Pour $p=2$ on retrouve le MILP précédent. Pour un objet de $\mathcal{C}^h$, seules $b^{h-1}$ et $b^h$ apparaissent dans ses
contraintes, donc dans le code on ne crée les variables que pour ces deux frontières (au plus $2n$ binaires par objet).
C'est ce programme que l'on résout avec Gurobi dans la partie suivante.


% ==================================================================
\section*{Jalon 1 : implémentation et validation}

\subsection*{Le programme}

Le MILP (cas à $p$ catégories) est codé en Python avec \texttt{gurobipy}, dans \texttt{src/mrsort.py}.
Le programme lit un learning set au format CSV (une ligne par objet : les $n$ performances puis la classe), en déduit
$n$ et $p$, normalise les données, résout le MILP et affiche $w$, $\lambda$, les frontières (dans l'échelle d'origine),
$\alpha^*$ et le taux de bonne classification sur le learning set.

\subsection*{Modèle appris sur Dataset0}

Dataset0 contient 60 objets, 5 critères à maximiser et 3 catégories. On obtient en moins d'une seconde un modèle
qui classe correctement les 60 objets ($\alpha^*=0.071$) :

\begin{center}
\begin{tabular}{lccccc}
\toprule
 & crit1 & crit2 & crit3 & crit4 & crit5 \\
\midrule
Poids $w_i$ & $1/7$ & $2/7$ & $1/7$ & $1/7$ & $2/7$ \\
Frontière $b^1$ (catégorie 1 / 2) & 34.5 & 35 & 33 & 34.5 & 33 \\
Frontière $b^2$ (catégorie 2 / 3) & 41.5 & 49 & 33 & 43.5 & 49 \\
\bottomrule
\end{tabular}
\qquad $\lambda = 0.643 \approx 4.5/7$
\end{center}

Les critères 2 et 5 comptent double. Pour franchir une frontière, il faut l'atteindre sur des critères qui totalisent
au moins 5 points sur 7 (par exemple crit2, crit5 et un troisième). Ce modèle n'est pas le seul possible (cf. question iii en annexe).
Sur chaque critère, la frontière peut être n'importe où entre deux valeurs voisines des données sans changer les affectations.
On l'a placée au milieu de cet intervalle pour qu'elle soit le plus loin possible des données.

\subsection*{Protocole de test}

Pour voir si le programme apprend bien, on part d'un modèle connu et on regarde si on le retrouve :
\begin{enumerate}
  \item On tire au hasard un modèle MR-Sort de référence $\mathcal{M}'$ : poids, $\lambda\in[0.5,1]$, frontières ordonnées.
  \item On tire $|L|$ objets aléatoires dans $[0,1]^4$ et on les classe avec $\mathcal{M}'$ : c'est le learning set $L$.
  \item On apprend $\mathcal{M}^*$ sur $L$ avec notre programme.
  \item On tire 1000 nouveaux objets $T$, jamais vus, et on les classe avec $\mathcal{M}'$ et $\mathcal{M}^*$.
        La \emph{qualité} est la proportion d'objets affectés à la même classe.
\end{enumerate}
On répète 20 fois et on fait la moyenne. On teste sur de nouveaux objets parce que sur $L$ le modèle appris est
correct à 100\,\% par construction.

\subsection*{Résultats}

\begin{center}
\begin{tabular}{ccccc}
\toprule
$p$ & $|L|$ & Qualité moyenne sur $T$ & Écart-type & Temps moyen \\
\midrule
2 & 20  & 88.5\,\% & 5.4\,\% & 0.05 s \\
2 & 50  & 96.4\,\% & 2.5\,\% & 0.18 s \\
2 & 100 & 98.4\,\% & 1.2\,\% & 0.49 s \\
3 & 50  & 91.5\,\% & 4.3\,\% & 0.31 s \\
3 & 100 & 96.6\,\% & 2.0\,\% & 0.74 s \\
\bottomrule
\end{tabular}
\end{center}

Dans les 100 essais, le learning set est bien classé à 100\,\%. La qualité sur $T$ augmente avec la taille de $L$,
car avec plus d'exemples, les frontières sont mieux encadrées et $\mathcal{M}^*$ se rapproche de $\mathcal{M}'$. Avec 3 catégories
il faut plus d'exemples, puisqu'il y a deux frontières à apprendre.


% ==================================================================
\newpage
\section*{Annexe : questions préliminaires (admission d'étudiants)}

% ------------------------------------------------------------------
\subsection*{Question i}

Somme des poids des cours validés ($\geq 12$), avec $w=(0.25,0.15,0.2,0.4)$, $\lambda=0.66$ :
Alice $0.15+0.2+0.4=0.75$ ($\A$), Bob $0.25+0.15+0.2=0.60$ ($\R$),
Charlie $0.15+0.4=0.55$ ($\R$), David $1$ ($\A$).
Bob ne peut pas améliorer sa situation. Avec 12 en physique, ce cours est déjà validé et une meilleure note
n'y change rien. Seule une meilleure note en histoire (10) pourrait le faire admettre.

% ------------------------------------------------------------------
\subsection*{Question ii}

Il faut que n'importe quels 3 cours atteignent $\lambda$ et que n'importe quels 2 cours restent en dessous. Le plus simple
est de donner le même poids aux 4 cours, $w_M=w_P=w_L=w_H=\frac14$, avec $\lambda\in\left]\frac12,\frac34\right]$, par exemple
$\lambda=0.75$. Alors 3 cours pèsent $\frac34\geq\lambda$ et 2 cours pèsent $\frac12<\lambda$. Ce n'est pas la seule solution,
par exemple $w=(0.3,0.25,0.25,0.2)$ avec $\lambda=0.6$ convient aussi.

% ------------------------------------------------------------------
\subsection*{Question iii}

Coalitions gagnantes (somme des poids $\geq \lambda$) dans les deux cas :

\begin{center}
\begin{tabular}{lcc}
\toprule
Coalition & $w = (0.49, 0.49, 0.02)$, $\lambda=0.5$ & $w = (\frac13,\frac13,\frac13)$, $\lambda=0.5$ \\
\midrule
$\{1\}$, $\{2\}$, $\{3\}$ & $0.49,\ 0.49,\ 0.02$ : perdantes & $0.33$ : perdantes \\
$\{1,2\}$ & $0.98$ : gagnante & $0.67$ : gagnante \\
$\{1,3\}$, $\{2,3\}$ & $0.51$ : gagnantes & $0.67$ : gagnantes \\
$\{1,2,3\}$ & $1$ : gagnante & $1$ : gagnante \\
\bottomrule
\end{tabular}
\end{center}

Dans les deux cas, un étudiant est admis si et seulement s'il valide \emph{au moins 2 cours sur 3}, donc
les deux jeux de paramètres définissent la même règle.

On en conclut que $(w,\lambda)$ n'est pas unique. Ce qui compte, ce sont les coalitions gagnantes. Les poids ne
donnent pas directement l'importance d'un cours (le cours 3 ne pèse que $0.02$ mais il est aussi décisif que les autres).
Pour juger un modèle appris, il faut donc regarder ses décisions plutôt que ses paramètres.

% ------------------------------------------------------------------
\subsection*{Question iv}

$b$ étant connu, on connaît pour chaque objet $x$ l'ensemble des critères qu'il valide, $C(x)=\{i\in\N : x_i\geq b_i\}$,
et son score $s(x)=\sum_{i\in C(x)} w_i$ est linéaire en $w$. Les variables sont les $w_i\geq 0$, $\lambda\in[0,1]$ et une marge $\alpha$.
\begin{align*}
\max\ & \alpha \\
\text{s.c.}\ & s(x) \geq \lambda + \alpha && \forall x\in A^* \\
& s(x) \leq \lambda - \alpha && \forall x\in R^* \\
& \textstyle\sum_{i\in\N} w_i = 1
\end{align*}
$\alpha$ est le plus petit écart entre un score et $\lambda$. En le maximisant, on éloigne le plus possible les acceptés et
les rejetés du seuil. Si $\alpha^*>0$, tous les objets sont bien classés, sinon aucun modèle ne colle parfaitement aux données.

\begin{center}
\begin{tikzpicture}[x=26cm,y=1cm,font=\small]
  \fill[gray!15] (0.60,-0.15) rectangle (0.75,0.9);
  \draw[->] (0.48,0) -- (1.03,0) node[right] {$s(x)$};
  \foreach \t in {0.5,0.6,0.7,0.8,0.9,1.0} \draw (\t,0.06) -- (\t,-0.06) node[below] {\footnotesize\t};
  \draw[dashed,thick] (0.675,-0.15) -- (0.675,1.15) node[above] {$\lambda$};
  \draw[<->] (0.60,0.7) -- node[above] {$\alpha$} (0.675,0.7);
  \draw[<->] (0.675,0.7) -- node[above] {$\alpha$} (0.75,0.7);
  \foreach \v/\n in {0.55/Charlie,0.60/Bob} \fill[red!75!black] (\v,0) circle (2.2pt) node[above=2pt,text=black] {\n};
  \foreach \v/\n in {0.75/Alice,1.0/David} \fill[green!50!black] (\v,0) circle (2.2pt) node[above=2pt,text=black] {\n};
\end{tikzpicture}

\footnotesize Scores des étudiants de la question i (rejetés en rouge, acceptés en vert).
\end{center}


\end{document}
